\section{Introduction}

Circular-economy monitoring uses changes in indicators to assess progress, compare national trajectories and judge the pace of material recirculation. These uses require more than comparable indicator levels: they require evidence that the temporal changes being interpreted are supported by the underlying observations. A sustained trend, a year-to-year change, a turning point and an acceleration summarize different features of a series and may have different evidential requirements. Indicator research has established that circularity measures must be interpreted in relation to their system boundaries and measurement purposes \parencite{Saidani2019,Moraga2019}. Their temporal information structure warrants the same attention.

The EU Circular Material Use Rate (CMUR) measures circular material use as a share of overall material use and is disseminated annually for Member States and the EU aggregate \parencite{Eurostat2018,EEAIndicator2026}. It combines waste-treatment statistics with material-flow accounts and trade information. Eurostat documents that the waste-treatment source is available every second year and that missing intervening years are estimated \parencite{EurostatMetadata2026}. An annual CMUR observation can therefore combine inputs with different observation frequencies. This structure matters when the rate is used to distinguish persistent movement from local temporal variation.

The source architecture is already part of CMUR research. \textcite{Christis2023} examine trends and the doubling target while explicitly recognizing even-year waste data and interpolation of odd years. \textcite{Nuss2024Nowcasting} likewise document the lower-frequency source in their work on nowcasting and outlooks. Most directly, \textcite{Haider2026} compare EU-aggregate annual and reporting-year trend estimates and their nominal precision as part of a forecasting assessment. These studies establish both the temporal uses of CMUR and prior awareness that annual dissemination incorporates estimated years. The remaining question concerns the temporal claims that this information can sustain across Member States.

Established temporal-disaggregation and benchmarking methods address the construction and consistency of higher-frequency statistics from lower-frequency information \parencite{ESS2018,Quilis2018}. Mixed-frequency composite-indicator research and empirical studies of environmental observation frequency further show why the information available at different time scales matters for estimation \parencite{Mattera2023,CortesGuzman2026}. For an already disseminated CMUR series, however, constructing annual values and assessing support for interpretations of those values are distinct tasks. Existing CMUR trend, target and forecasting analyses do not establish a Member-State-wide comparison of medium-term direction and precision with annual-change, turning and acceleration support under a common provenance-conditioned design.

We therefore ask: \emph{Which temporal claims remain supportable for the EU Circular Material Use Rate when source provenance and the lower observation frequency of the waste component are respected?} We study all 27 Member States and EU27\_2020 separately, comparing the released annual series with released CMUR restricted to verified admissible source-anchor years over the fixed 2010--2022 medium-term window. Identical two-sided 95\% Student-$t$ ordinary least-squares (OLS) rules distinguish supported trend direction from unresolved direction. Local temporal claims through 2022 are assessed with trajectories conditioned on the observed anchors.

The contribution is a CMUR-specific, Member-State-level map of a scale-dependent temporal-support boundary. Medium-term direction remains comparatively stable where both analyses resolve it, but restricting inference to verified anchors widens model-based uncertainty. Most annual-change, turning and acceleration comparisons remain unresolved under two paired residual-amplitude sensitivities and across the tested decision-threshold range. This comparison extends the EU-level reporting-year evidence by separating support for long-horizon direction from support for local temporal features across the full Member-State panel.

The evidence proceeds from medium-term direction and interval widths to local claim support. The short-window assessment retains the empirical four-year pseudo-gap residual amplitudes as its primary stress test and examines a secondary Brownian-motivated amplitude of $1/\sqrt{2}$ using the same sampled residuals. Support is then examined across the pre-specified 0.80--0.99 threshold grid. These scenarios describe sensitivity to unobserved intervening-year errors rather than identify their true two-year distribution. The resulting boundary concerns CMUR under this source-aware design: an unresolved comparison indicates insufficient support under the stated criteria, rather than absence of real change.
